\let\negmedspace\undefined
\let\negthickspace\undefined
\documentclass[journal]{article}
\usepackage[a5paper, margin=10mm, onecolumn]{geometry}
\usepackage{lmodern} % Ensure lmodern is loaded for pdflatex

\setlength{\headheight}{1cm} % Set the height of the header box
\setlength{\headsep}{0mm}     % Set the distance between the header box and the top of the text

\usepackage{gvv-book}
\usepackage{gvv}
\usepackage{cite}
\usepackage{textcomp}
\usepackage{amsmath,amssymb,amsfonts,amsthm}
\usepackage{algorithmic}
\usepackage{graphicx}
\graphicspath{{./figs/}}
\usepackage{textcomp}
\usepackage{xcolor}
\usepackage{txfonts}
\usepackage{listings}
\usepackage{enumitem}
\usepackage{mathtools}
\usepackage{gensymb}
\usepackage{comment}
\usepackage[breaklinks=true]{hyperref}
\usepackage{tkz-euclide} 
\usepackage{listings}
\usepackage{gvv}                                        
\def\inputGnumericTable{}                                 
\usepackage[latin1]{inputenc}                                
\usepackage{color}                                            
\usepackage{array}                                            
\usepackage{longtable}                                       
\usepackage{calc}                                             
\usepackage{multirow}                                         
\usepackage{hhline}                                           
\usepackage{ifthen}                                           
\usepackage{lscape}
\usepackage{circuitikz}
\tikzstyle{block} = [rectangle, draw, fill=blue!20, 
text width=4em, text centered, rounded corners, minimum height=3em]
\tikzstyle{sum} = [draw, fill=blue!10, circle, minimum size=1cm, node distance=1.5cm]
\tikzstyle{input} = [coordinate]
\tikzstyle{output} = [coordinate]


\begin{document}
	
	\bibliographystyle{IEEEtran}
	\vspace{3cm}
	
\title{9.2.7}
\author{EE25BTECH11047 - RAVULA SHASHANK REDDY}
\maketitle
\hrulefill
\bigskip 

\renewcommand{\thetable}{\theenumi}
\setlength{\intextsep}{10pt}

\textbf{Question:} \\

Find the area of the region bounded by the parabola $y^2 = x$ and the line $2y = x$.\\

\textbf{Solution:}
\begin{align}
\vec{x}^T \vec{V} \vec{x} + 2\vec{u}^T\vec{x} + f = 0 \\
\vec{V} = \myvec{0 & 0 \\ 0 & 1},\quad 
\vec{u} = \myvec{-\tfrac{1}{2} \\ 0},\quad f=0 \\
\end{align}
\begin{align}
\vec{x} &= \vec{h} + \kappa \vec{m}, \quad \kappa \in \mathbb{R} \\
\vec{h} &= \myvec{0 \\ 0}, \quad 
\vec{m} = \myvec{2 \\ 1} \\
\vec{x}_i &= \vec{h} + \kappa_i \vec{m} 
\end{align}

Substituting eq-(6) in eq-(1) and Solving the Quadratic 
\begin{align}
\kappa_i = \frac{1}{\vec{m}^T \vec{V} \vec{m}}
\brak{-\vec{m}^T(\vec{V}\vec{h}+\vec{u})
\pm \sqrt{(\vec{m}^T(\vec{V}\vec{h}+\vec{u}))^2 - g(\vec{h})(\vec{m}^T \vec{V} \vec{m})}} 
\end{align}
\begin{align}
g(\vec{h}) &= \vec{h}^T \vec{V} \vec{h} + 2\vec{u}^T \vec{h} + f\\
\vec{V}\vec{m} &= \myvec{0 \\ 1},\quad 
\vec{m}^T \vec{V} \vec{m} = 1 \\
\vec{V}\vec{h}+\vec{u} &= \myvec{-\tfrac{1}{2} \\ 0},\quad
\vec{m}^T(\vec{V}\vec{h}+\vec{u}) = -1 \\
g(\vec{h}) &= 0\\
\kappa_i &= -(-1) \pm \sqrt{(-1)^2 - 0} 
= 1 \pm 1 \\
\implies \kappa_1 &= 0,\;\; \kappa_2 = 2
\end{align}

\begin{align}
\vec{x}_1 &= \vec{h} + 0\vec{m} = \myvec{0 \\ 0},\quad
\vec{x}_2 = \vec{h} + 2\vec{m} = \myvec{4 \\ 2} \\
\text{Area} &= \int_{0}^{2} \Big(2y - y^2\Big)\,dy 
= \brak{y^2 - \tfrac{y^3}{3}}_0^2 
= 4 - \tfrac{8}{3} 
= \tfrac{4}{3}
\end{align}
\newpage
\begin{figure}
    \centering
    \includegraphics[width=1.0\linewidth]{figs/Fig_1.png}
    \caption{}
    \label{fig:placeholder}
\end{figure}
\end{document}